\begin{proof}[Proof of \cref{obj:prop:dorn-a3-free-corollary}]
\begin{enumerate}
\item \emph{The source mean is bounded throughout the cube.}
Fix \(\gamma\in\Gamma\) and a source tuple in \(\mathcal D_\gamma\). Write
\[
X^{\mathrm D}:=X,
\qquad
P_X^{\mathrm D}:=\mathcal L_P(X^{\mathrm D})
=2^{-d}\operatorname{Leb}\big|_{[-1,1]^d},
\qquad
\mu_{1,P}:=\mu_P,
\]
and let \(\mu_{0,P}\) be its selected control mean. Continuity of \(\mu_P\) on the compact source cube gives
\[
\int |\mu_P(z)|\,P_X^{\mathrm D}(dz)<\infty.
\]
For \(P_X^{\mathrm D}\)-almost every \(z\), the treated conditional distribution is
\(N(\mu_P(z),\sigma_{\min}^2)\). Its mean identity, the absolute-integral inequality, and Jensen's inequality give
\[
|\mu_P(z)|
\le
\int |y|\,N(\mu_P(z),\sigma_{\min}^2)(dy),
\]
and
\[
|\mu_P(z)|^q
\le
\left(\int |y|\,N(\mu_P(z),\sigma_{\min}^2)(dy)\right)^q
\le
\int |y|^q\,N(\mu_P(z),\sigma_{\min}^2)(dy)
\le M^q.
\]
Since \(q>0\) and \(M>0\), this implies \(|\mu_P|\le M\) almost everywhere.

To extend this inequality to the boundary as well, define
\[
v_P(z):=\max\{|\mu_P(z)|-M,0\},
\qquad z\in\mathbb R^d.
\]
The function \(v_P\) is continuous on the source cube and vanishes Lebesgue-almost everywhere there. Moreover,
\[
[-1,1]^d
=
\overline{\operatorname{int}([-1,1]^d)}.
\]
Continuity first forces \(v_P\) to vanish on the interior, since a positive value would give a neighborhood of positive Lebesgue measure on which it remains positive. Taking limits from the interior then gives
\[
v_P(z)=0,
\qquad
|\mu_P(z)|\le M
\qquad(z\in[-1,1]^d).
\]
% lean: CausalSmith.Stat.WeakOverlap.dorn_A3_free_corollary

\item \emph{Completion and comparison of the common radii.}
By \cref{obj:lem:fixed-cube-holder-completion}, there is a constant \(K_{d,\beta}>0\) such that, for every nonnegative amplitude bound \(M'\), nonnegative seminorm bound \(L_0'\), and function satisfying the corresponding source regularity and bounds,
\[
u\circ T
\in
\mathcal H^\beta\!\left(K_{d,\beta}(M'+L_0')\right).
\]
For a function admitting a nonnegative Hölder radius, write its set of admissible radii as
\[
\mathscr L_\beta(g)
:=
\{L_{\mathrm{adm}}\in[0,\infty):
g\in\mathcal H^\beta(L_{\mathrm{adm}})\},
\qquad
\|g\|_{\mathcal H^\beta}
=
\inf\mathscr L_\beta(g).
\]
Consequently,
\[
\|u\circ T\|_{\mathcal H^\beta}
\le K_{d,\beta}(M'+L_0').
\]
In particular, the two suprema in the proposition are bounded above:
\[
C_{d,\beta}\le K_{d,\beta},
\qquad
L_*\le K_{d,\beta}(M+L_0).
\]

Every admissible radius bounds the function itself. Taking the infimum of those bounds gives
\[
|g(x)|\le \|g\|_{\mathcal H^\beta},
\qquad x\in[0,1]^d.
\]
Constant functions satisfy the source seminorm restriction because their derivative differences vanish. Applying the preceding inequality to the constant functions \(1\) and \(M\), which admit radii by the completion result, yields
\[
1\le \|1\|_{\mathcal H^\beta},
\qquad
M\le \|M\|_{\mathcal H^\beta}.
\]
The constant function \(1\), with amplitude and seminorm bounds both equal to \(1\), is admissible in the definition of \(C_{d,\beta}\). The constant function \(M\) is admissible in the definition of \(L_*\). Hence
\[
0<\frac12\le C_{d,\beta}\le K_{d,\beta}<\infty,
\qquad
0<M\le L_*.
\]
Finally, every function contributing to \(L_*\) also contributes, after normalization, to \(C_{d,\beta}\):
\[
\frac{\|u\circ T\|_{\mathcal H^\beta}}{M+L_0}
\le C_{d,\beta}.
\]
Taking the supremum over these functions proves
\[
0<L_*\le C_{d,\beta}(M+L_0)<\infty.
\]
% lean: CausalSmith.Stat.WeakOverlap.dorn_A3_free_corollary

\item \emph{Gaussian source witnesses for the local-condition assertion.}
Suppose \(d\ge2\), and fix \(\gamma\in\Gamma\). Then \(\gamma>1\). Choose
\[
\delta_{\mathrm{strip}}:=\frac{d}{\gamma-1}>0,
\qquad
a:=1+\frac{\delta_{\mathrm{strip}}}{2},
\]
so that
\[
1<a<1+\frac{d}{\gamma-1}.
\]
For the fixed moment exponent, set
\[
I_q:=\int |y|^q\,N(0,1)(dy),
\qquad
M':=\max\{1,|I_q|\},
\qquad
C':=\left((3/4)^{\gamma-1}\right)^{-1},
\qquad
\sigma_{\min}':=1,
\qquad
L_0':=1.
\]
Gaussian moments are finite, so \(M'\) is positive and finite. Since \(M'\ge1\) and \(q>1\),
\[
I_q\le |I_q|\le M'\le (M')^q.
\]
All four chosen source constants are positive and finite.

Use the center, radial function, strip, and propensity from
\cref{obj:def:thin-strip}. For the ambient arguments needed below, the same coordinate formulas are
\[
x^\circ=(1/2,\ldots,1/2),
\qquad
R(x)=\|x-x^\circ\|_\infty
\quad(x\in\mathbb R^d),
\]
\[
F_R(t)=
\begin{cases}
0,&t<0,\\
\min\{1,(2t)^d\},&t\ge0,
\end{cases}
\qquad
\widetilde S:=
\{x\in\mathbb R^d:
|x_2-1/2|\le |x_1-1/2|^a\}.
\]
Here \(\widetilde S\cap[0,1]^d=S\). Define the capped propensity on all of \(\mathbb R^d\) by
\[
e_{\mathrm{cap}}(x)
:=
\min\left\{\frac34,\,
\min\left\{1,\,
F_R(R(x))^{1/(\gamma-1)}
+\frac14\mathbf1_{\widetilde S}(x)\right\}\right\}.
\]
Thus, on the unit cube,
\[
e_{\mathrm{cap}}(x)=\min\{3/4,e_\star(x)\}.
\]

Construct random variables with conditional product distribution
\[
X_{\mathrm{strip}}\sim\operatorname{Unif}([0,1]^d),
\]
\[
\mathcal L\!\left(
A_{\mathrm{strip}},Y_{\mathrm{strip}}(0),Y_{\mathrm{strip}}(1)
\mid X_{\mathrm{strip}}=x
\right)
=
\operatorname{Bernoulli}(e_{\mathrm{cap}}(x))
\otimes N(0,1)\otimes N(0,1).
\]
Set
\[
Y_{\mathrm{strip}}
:=
A_{\mathrm{strip}}Y_{\mathrm{strip}}(1)
+(1-A_{\mathrm{strip}})Y_{\mathrm{strip}}(0),
\qquad
X_{\mathrm G}^{\mathrm D}:=T(X_{\mathrm{strip}}),
\]
\[
P_{\mathrm G}:=
\mathcal L(X_{\mathrm G}^{\mathrm D},A_{\mathrm{strip}},Y_{\mathrm{strip}}),
\qquad
\mu_{1,\mathrm G}(z)=\mu_{0,\mathrm G}(z):=0,
\qquad
e_{\mathrm G}(z):=e_{\mathrm{cap}}(T^{-1}z)
\quad(z\in\mathbb R^d).
\]
These measurable product kernels define a probability law. The affine change of variables and the conditional product distribution give
\[
X_{\mathrm G}^{\mathrm D}\sim\operatorname{Unif}([-1,1]^d),
\qquad
e_{\mathrm G}(X_{\mathrm G}^{\mathrm D})
=
P_{\mathrm G}(A_{\mathrm{strip}}=1\mid X_{\mathrm G}^{\mathrm D}),
\]
and
\[
Y_{\mathrm{strip}}
\mid X_{\mathrm G}^{\mathrm D},A_{\mathrm{strip}}=b
\sim N(0,1),
\qquad b\in\{0,1\}.
\]
Since the radial term is positive away from \(x^\circ\), a set of full covariate probability,
\[
0<e_{\mathrm G}(X_{\mathrm G}^{\mathrm D})\le\frac34<1
\quad\text{almost surely}.
\]
The zero means satisfy the source smoothness restriction, and
\[
\operatorname{Var}_{P_{\mathrm G}}(\mu_{1,\mathrm G}(X_{\mathrm G}^{\mathrm D}))=0\le M',
\qquad
\operatorname{Var}_{P_{\mathrm G}}(Y_{\mathrm{strip}}\mid
X_{\mathrm G}^{\mathrm D},A_{\mathrm{strip}})=1=(\sigma_{\min}')^2.
\]
Together with the bound on \(I_q\), these establish the moment, variance, Gaussian, and smoothness restrictions.

The uniform-covariate propensity bound in
\cref{obj:prop:strict-enlargement} supplies
\[
\Pr\{e_\star(X_{\mathrm{strip}})\le t\}\le t^{\gamma-1},
\qquad 0\le t\le1.
\]
For \(0\le t<3/4\), capping leaves this sublevel event unchanged. Since \(C'\ge1\),
\[
\Pr\{e_{\mathrm{cap}}(X_{\mathrm{strip}})\le t\}
\le t^{\gamma-1}\le C't^{\gamma-1}.
\]
For \(3/4\le t\le1\), monotonicity of the positive power gives
\[
\Pr\{e_{\mathrm{cap}}(X_{\mathrm{strip}})\le t\}
\le1
=
C'(3/4)^{\gamma-1}
\le C't^{\gamma-1}.
\]
The affine transport therefore yields
\[
P_{\mathrm G}\{e_{\mathrm G}(X_{\mathrm G}^{\mathrm D})\le t\}
\le C't^{\gamma-1},
\qquad 0\le t\le1.
\]
Thus the tuple satisfies every retained source restriction with the chosen constants.

It remains to establish the failure of the numerical clause in
\cref{obj:def:dorn-a3} for this capped representative. Define the closed centered Euclidean ball by
\[
\mathcal B_h:=
\{x\in\mathbb R^d:\|x-x^\circ\|_2\le h\},
\qquad h>0.
\]
For \(0<h\le1/2\), the ball lies in the unit cube. On
\(S\cap\mathcal B_h\), the second coordinate has width at most \(2h^a\), and each other coordinate has width at most \(2h\). Conversely,
\[
\prod_{i=1}^d[1/2-h/d,1/2+h/d]\subseteq\mathcal B_h,
\]
because the squared Euclidean distance of any point in this box is at most
\(d(h/d)^2\le h^2\). These box bounds give
\[
\Pr\{X_{\mathrm{strip}}\in S\cap\mathcal B_h\}
\le (2h^a)(2h)^{d-1},
\]
\[
\Pr\{X_{\mathrm{strip}}\in\mathcal B_h\}
\ge (2h/d)^d>0,
\]
and hence
\[
\frac{\Pr\{X_{\mathrm{strip}}\in S\cap\mathcal B_h\}}
{\Pr\{X_{\mathrm{strip}}\in\mathcal B_h\}}
\le
\frac{(2h^a)(2h)^{d-1}}{(2h/d)^d}
=
d^d h^{a-1}.
\]

Consider any proposed constants \(\rho,\nu,h_0>0\), and set
\[
\rho_{\mathrm{cap}}:=\min\{\rho,2\}>0.
\]
Both \(h^{a-1}\) and \(((2h)^d)^{1/(\gamma-1)}\) tend to zero as \(h\downarrow0\). Choose \(h\) with
\[
0<h\le\min\{h_0,1/2\},
\qquad
d^d h^{a-1}\le\nu,
\qquad
((2h)^d)^{1/(\gamma-1)}
<
\frac{\rho_{\mathrm{cap}}}{4}.
\]
For every \(x\in\mathcal B_h\), \(R(x)\le h\), so
\[
F_R(R(x))^{1/(\gamma-1)}
\le ((2h)^d)^{1/(\gamma-1)}
<
\frac{\rho_{\mathrm{cap}}}{4}
\le\frac12.
\]
Consequently,
\[
e_\star(x)
\le F_R(R(x))^{1/(\gamma-1)}+\frac14
<\frac34,
\qquad
e_{\mathrm{cap}}(x)=e_\star(x)
\quad(x\in\mathcal B_h).
\]
At the center the propensity equals \(1/4\), giving the pointwise bound
\[
e_{\mathrm{cap}}(x^\circ)=\frac14,
\qquad
\sup_{x\in\mathcal B_h\cap[0,1]^d}e_{\mathrm{cap}}(x)\ge\frac14.
\]
For \(x\in\mathcal B_h\setminus S\), the strip term vanishes, and
\[
e_{\mathrm{cap}}(x)
=
F_R(R(x))^{1/(\gamma-1)}
<\frac{\rho}{4}.
\]
It follows that
\[
\left\{
x\in\mathcal B_h:
e_{\mathrm{cap}}(x)\ge
\rho\sup_{x'\in\mathcal B_h\cap[0,1]^d}e_{\mathrm{cap}}(x')
\right\}
\subseteq S\cap\mathcal B_h.
\]
The resulting local mass ratio is therefore at most
\(d^d h^{a-1}\le\nu\), contradicting the required strict inequality.

Finally, the affine identities
\[
\|T(x)-T(x^\circ)\|_2=2\|x-x^\circ\|_2,
\qquad
T(x^\circ)=\mathbf0,
\qquad
e_{\mathrm G}(T(x))=e_{\mathrm{cap}}(x)
\]
preserve the pointwise suprema and the numerator and denominator probabilities in these ratios. A source-cube local condition with radius bound \(h_0\) would thus give a unit-cube condition with radius bound \(h_0/2\), which the preceding argument contradicts. Hence
\[
\neg\,\mathsf A_3^{\mathrm{Dorn}}
\bigl(\{P_{\mathrm G}\},(e_{\mathrm G})\bigr).
\]
The construction selects its constants separately for the chosen \(\gamma\).
% lean: CausalSmith.Stat.WeakOverlap.gaussianThinStripOriginal_not_dornA3

\item \emph{Transport and causal completion of each source tuple.}
Return to an arbitrary tuple in \(\mathcal D_\gamma\), with \(\gamma\in\Gamma\). Define globally measurable versions by
\[
\mu_{b,P}^{\mathrm{meas}}(z):=
\begin{cases}
\mu_{b,P}(z),&z\in[-1,1]^d,\\
0,&z\notin[-1,1]^d,
\end{cases}
\qquad
b\in\{0,1\},\ z\in\mathbb R^d,
\]
\[
e_P^{\mathrm{cl}}(z):=
\max\{0,\min\{1,e_P(z)\}\},
\qquad z\in\mathbb R^d.
\]
Continuity on the closed source cube makes the mean extensions measurable. The clamped propensity is measurable and takes values in \([0,1]\). Since the source covariate is supported on the cube and the selected propensity lies in \((0,1)\) almost surely,
\[
\mu_{b,P}^{\mathrm{meas}}=\mu_{b,P},
\qquad
e_P^{\mathrm{cl}}=e_P
\quad P_X^{\mathrm D}\text{-almost everywhere}.
\]

Construct a completion law \(\Pi_P\) by retaining covariate marginal
\(P_X^{\mathrm D}\) and using the conditional product distribution
\[
\mathcal L_{\Pi_P}(A,Y(0),Y(1)\mid X^{\mathrm D}=z)
=
\operatorname{Bernoulli}(e_P^{\mathrm{cl}}(z))
\otimes
N(\mu_{0,P}^{\mathrm{meas}}(z),\sigma_{\min}^2)
\otimes
N(\mu_{1,P}^{\mathrm{meas}}(z),\sigma_{\min}^2),
\]
with
\[
Y:=AY(1)+(1-A)Y(0).
\]
All factors are probability kernels, so this defines a probability law.

To identify its observed marginal, take any bounded measurable function
\[
\varphi:\mathbb R^d\times\{0,1\}\times\mathbb R\longrightarrow\mathbb R.
\]
Conditioning first on the covariate and treatment and then on the covariate alone gives
\[
\begin{aligned}
\int \varphi\,dP
={}&
\int\Bigg[
(1-e_P^{\mathrm{cl}}(z))
\int\varphi(z,0,y)\,
N(\mu_{0,P}^{\mathrm{meas}}(z),\sigma_{\min}^2)(dy)
\\
&\hspace{34mm}
+
e_P^{\mathrm{cl}}(z)
\int\varphi(z,1,y)\,
N(\mu_{1,P}^{\mathrm{meas}}(z),\sigma_{\min}^2)(dy)
\Bigg]P_X^{\mathrm D}(dz)
\\
={}&
\mathbb E_{\Pi_P}\varphi(X^{\mathrm D},A,Y).
\end{aligned}
\]
Here the selected propensity and both Gaussian arm kernels identify the two conditional distributions. Thus
\[
\mathcal L_{\Pi_P}(X^{\mathrm D},A,Y)=P.
\]

Define the transported covariate, completion, and observed law by
\[
X^T:=T^{-1}(X^{\mathrm D}),
\]
\[
\Pi_P^T:=
\mathcal L_{\Pi_P}(X^T,A,Y(0),Y(1),Y),
\qquad
Q_P:=
\mathcal L_P(T^{-1}(X^{\mathrm D}),A,Y).
\]
The observed marginal of \(\Pi_P^T\) is exactly \(Q_P\). The source tail inequality at \(t=1\) also gives
\[
1=P\{e_P(X^{\mathrm D})\le1\}\le C.
\]
Together with \(B_*>0\), \(L_*>0\), and \(\gamma>1\), this verifies the required parameter domain.

The affine change of variables has Jacobian \(2^d\), and hence
\[
X^T\sim\operatorname{Unif}([0,1]^d),
\qquad
f_{Q_P}(x)=1
\quad\text{Lebesgue-almost everywhere on }[0,1]^d.
\]
The transported selected propensity is
\[
e_P^T(x):=e_P(T(x)),
\qquad x\in\mathbb R^d,
\]
and satisfies
\[
Q_P\{e_P^T(X^T)\le t\}
=
P\{e_P(X^{\mathrm D})\le t\}
\le Ct^{\gamma-1},
\qquad 0\le t\le1.
\]
The treated conditional kernel under \(Q_P\) is
\(N(\mu_P(T(x)),\sigma_{\min}^2)\) almost everywhere. The displayed bound \(|\mu_P(z)|\le M\) on the source cube gives
\[
|\mu_P(T(x))|\le M\le B_*,
\qquad x\in[0,1]^d.
\]
Translation of a Gaussian by its mean produces a centered Gaussian. Its moment-generating function therefore gives, for every \(\lambda\in\mathbb R\),
\[
\mathbb E_{Q_P}\!\left[
\exp\{\lambda(Y-\mu_P(T(X^T)))\}
\mid X^T,A=1
\right]
=
\exp\!\left(\frac{\lambda^2\sigma_{\min}^2}{2}\right)
\le
\exp\!\left(\frac{\lambda^2B_*^2}{2}\right).
\]
These inequalities verify \cref{obj:ass:bounded-potential-outcomes}.

The treated potential outcome is integrable, by integrability of the source mean and the Gaussian residual. Its Gaussian conditional distribution yields
\[
\mathbb E_{\Pi_P}[Y(1)\mid X^{\mathrm D}]
=
\mu_{1,P}^{\mathrm{meas}}(X^{\mathrm D}).
\]
Since \(T\) is a measurable bijection with measurable inverse, conditioning on \(X^T\) gives
\[
\mathbb E_{\Pi_P^T}[Y(1)\mid X^T]
=
\mu_{1,P}^{\mathrm{meas}}(T(X^T))
=
\mu_P(T(X^T))
\quad\text{almost surely}.
\]

We next verify smoothness at the precise radius \(L_*\). Completion supplies at least one admissible radius for \(\mu_P\circ T\), and its defining supremum gives
\[
\|\mu_P\circ T\|_{\mathcal H^\beta}\le L_*.
\]
For completeness, the infimum defining the full radius is attained whenever its admissible-radius set is nonempty. Indeed, regularity follows from any one admissible radius, and every derivative bound satisfies
\[
|D^\alpha g(x)|
\le
\inf\mathscr L_\beta(g)
=
\|g\|_{\mathcal H^\beta},
\qquad |\alpha|\le m.
\]
For distinct \(x,x'\), every admissible radius bounds the quotient, so
\[
\frac{|D^\alpha g(x)-D^\alpha g(x')|}
{\|x-x'\|_\infty^{\beta-m}}
\le
\inf\mathscr L_\beta(g)
=
\|g\|_{\mathcal H^\beta},
\qquad |\alpha|=m.
\]
For \(x=x'\), both sides of the corresponding modulus inequality have zero left-hand side and zero distance factor; it follows directly from any admissible radius. Thus
\[
g\in\mathcal H^\beta(\|g\|_{\mathcal H^\beta}).
\]
Increasing a radius preserves its defining inequalities. Applying this to the transported response proves
\[
\mu_P\circ T\in\mathcal H^\beta(L_*).
\]
The measurable mean extension agrees with the selected source mean throughout the source cube, so its transported version agrees with \(\mu_P\circ T\) throughout the unit cube. Their intrinsic derivative traces and modulus restrictions therefore coincide. Since \(X^T\) is supported there, this also preserves the conditional-mean identity.

Finally, the product construction and its affine transport give
\[
Y=AY(1)+(1-A)Y(0)
\quad\text{almost surely},
\qquad
(Y(0),Y(1))\perp A\mid X^T.
\]
These are \cref{obj:ass:consistency,obj:ass:exchangeability}. Together with the density, tail, response-bound, and smoothness calculations, they verify every restriction in \cref{obj:def:model}. Consequently,
\[
Q_P\in\mathcal P_\gamma
\quad\text{with}\quad
(B,L,C,c_f)=(B_*,L_*,C,1),
\]
with the exhibited causal completion \(\Pi_P^T\).
% lean: CausalSmith.Stat.WeakOverlap.dorn_A3_free_corollary

\item \emph{A common expected-supremum bound.}
Define the enlarged tail constant by
\[
C_{\mathrm{up}}:=\max\{C,1\}.
\]
Increasing the tail constant preserves the tail restriction, since
\[
Ct^{\gamma-1}\le C_{\mathrm{up}}t^{\gamma-1},
\qquad 0\le t\le1.
\]
Thus every transported tuple belongs to the model with parameters
\((d,\beta,B_*,L_*,C_{\mathrm{up}},1)\). These parameters satisfy the hypotheses of \cref{obj:thm:adaptive-minimax}.

To state the sampling identity explicitly, define, for \(n\in\mathbb N\),
\[
\Omega_n:=
(\mathbb R^d\times\{0,1\}\times\mathbb R)^n,
\]
\[
\mathcal T_n\bigl((z_i,a_i,y_i)_{i=1}^n\bigr)
:=
(T^{-1}z_i,a_i,y_i)_{i=1}^n,
\qquad
\mathcal T_n:\Omega_n\to\Omega_n.
\]
For \(n=0\), this is the identity on the singleton empty sample. For every source tuple, define the transported-sample loss by
\[
Z_{n,P}(\omega)
:=
\sup_{x\in[0,1]^d}
|\widehat\mu_n(\omega,x)-\mu_P(T(x))|,
\qquad \omega\in\Omega_n,
\]
where \(\widehat\mu_n\) is the estimator of \cref{obj:def:estimator} with parameters \(\beta,B_*\).

By \cref{obj:lem:causal-identification}, the continuous treated response associated with \(Q_P\) agrees throughout the cube with the response of the displayed completion:
\[
\mu_{1,Q_P}(x)=\mu_P(T(x)),
\qquad x\in[0,1]^d.
\]
Hence the estimator risk supplied by \cref{obj:thm:adaptive-minimax} is exactly
\[
\mathbb E_{Q_P^{\otimes n}}Z_{n,P}
=
\mathbb E_{P^{\otimes n}}
\bigl[Z_{n,P}(\mathcal T_n(\mathcal D_n))\bigr].
\]
That result supplies \(K>0\) and a common integer \(n_0\ge1\) such that
\[
\mathbb E_{Q_P^{\otimes n}}Z_{n,P}
\le K r_{\gamma,n},
\qquad
n\ge n_0,\quad \gamma\in\Gamma,\quad P\in\mathcal D_\gamma.
\]
The constants are common because the model parameters and adaptation interval are fixed.
% lean: CausalSmith.Stat.WeakOverlap.adaptive_minimax_rate

\item \emph{Restriction to source families.}
The source setting permits an arbitrary nonempty family whose members satisfy the retained restrictions with common constants
\citep[Section 2, Assumptions 1--2, pp.~3--7]{Dorn2026GlobalOverlap}.
Therefore every such family satisfies
\[
\mathfrak P\ne\varnothing,
\qquad
\mathfrak P\subseteq\mathcal D_\gamma.
\]
Taking the supremum of the preceding memberwise estimate over this family gives, with the same \(K,n_0\),
\[
\sup_{P\in\mathfrak P}
\mathbb E_{P^{\otimes n}}
\bigl[Z_{n,P}(\mathcal T_n(\mathcal D_n))\bigr]
\le K r_{\gamma,n},
\qquad n\ge n_0.
\]
This is the asserted family expected-supremum bound.
% lean: CausalSmith.Stat.WeakOverlap.dorn_A3_free_corollary

\item \emph{The probability bound on the source cube.}
Fix \(\varepsilon>0\), and choose
\[
c(\varepsilon):=\frac K\varepsilon>0,
\qquad
n_{0,\varepsilon}:=\max\{n_0,1\}.
\]
For every \(n\ge n_{0,\varepsilon}\), the oracle scale is positive:
\[
r_{\gamma,n}
=
n^{-\beta/(2\beta+d\gamma/(\gamma-1))}>0.
\]
For each partition cell \(Q\) at each fixed dyadic mesh, choose a countable subset dense in the relative topology:
\[
D_Q\subseteq Q,
\qquad
\overline{D_Q}^{\,Q}=Q.
\]
On a sample with positive selected mesh, the fitted function restricted to each selected cell is a clipped polynomial. Together with continuity of the transported response, this gives
\[
\sup_{x\in Q}
|\widehat\mu_n(\omega,x)-\mu_P(T(x))|
=
\sup_{x\in D_Q}
|\widehat\mu_n(\omega,x)-\mu_P(T(x))|,
\qquad Q\in\mathcal Q_{\widehat h(\omega)}.
\]
For each fixed mesh and cell, the fitted value at a fixed point is sample-measurable, so the right-hand side is a countable supremum of measurable functions. Taking the maximum over the finitely many cells preserves measurability. The mesh selector is measurable and takes values in a countable set of dyadic meshes together with zero, so selecting among these measurable losses also preserves measurability. On the zero-mesh branch,
\[
Z_{n,P}(\omega)
=
\sup_{x\in[0,1]^d}|\mu_P(T(x))|,
\qquad \widehat h(\omega)=0,
\]
which is deterministic. Thus \(Z_{n,P}\) is measurable.

For every transported sample \(\omega\) and every source-cube point \(z\),
\[
\left|
\widehat\mu_n(\omega,T^{-1}z)-\mu_P(z)
\right|
=
\left|
\widehat\mu_n(\omega,T^{-1}z)
-\mu_P(T(T^{-1}z))
\right|
\le Z_{n,P}(\omega).
\]
The source covariate law is supported on that cube, so
\[
\|\widehat\mu_n(\omega,\cdot)\circ T^{-1}-\mu_P\|_{L^\infty(P)}
\le Z_{n,P}(\omega).
\]
Markov's inequality and the common expectation bound now give
\[
\begin{aligned}
&Q_P^{\otimes n}\!\left\{
\|\widehat\mu_n\circ T^{-1}-\mu_P\|_{L^\infty(P)}
>c(\varepsilon)r_{\gamma,n}
\right\}
\\
&\qquad\le
Q_P^{\otimes n}\{Z_{n,P}>c(\varepsilon)r_{\gamma,n}\}
\le
\frac{\mathbb E_{Q_P^{\otimes n}}Z_{n,P}}
{(K/\varepsilon)r_{\gamma,n}}
\le\varepsilon.
\end{aligned}
\]
The sampling identity for \(\mathcal T_n\) expresses the same probability under \(P^{\otimes n}\) when the estimator is computed from the transported sample. Every bound is simultaneous over \(\gamma\in\Gamma\) and source tuples in \(\mathcal D_\gamma\).
% lean: CausalSmith.Stat.WeakOverlap.measurableLoss_probability_le

\item \emph{Families with a common regression curve.}
Let \(\mathfrak P\) be the stated nonempty family with common treated curve \(\mu^\circ\). For any \(n\in\mathbb N\) and evaluation point \(x_0\), define the constant sample estimator
\[
\widehat\theta_{\mathrm{known},n}(\omega):=\mu^\circ(x_0),
\qquad \omega\in\Omega_n.
\]
It is measurable, including when \(n=0\). For every family member and every sample,
\[
|\widehat\theta_{\mathrm{known},n}(\omega)-\mu_P(x_0)|=0.
\]
Since \(cr>0\), its miss event is empty. Consequently,
\[
\begin{aligned}
0
&\le
\inf_{\substack{\widehat\theta:\Omega_n\to\mathbb R\\
\widehat\theta\ \mathrm{measurable}}}
\sup_{P\in\mathfrak P}
P^{\otimes n}\{|\widehat\theta-\mu_P(x_0)|>cr\}
\\
&\le
\sup_{P\in\mathfrak P}
P^{\otimes n}
\{|\widehat\theta_{\mathrm{known},n}-\mu_P(x_0)|>cr\}
=0.
\end{aligned}
\]
This proves the asserted equality in the extended nonnegative reals.
% lean: CausalSmith.Stat.WeakOverlap.knownRegression_minimax_miss_zero

\item \emph{The causal-class pointwise lower bound.}
Apply the pointwise lower conclusion, supported by the two-valued potential-outcome subclass, of
\cref{obj:thm:adaptive-minimax} with the fixed parameters
\[
(B,L,C,c_f)=(B_*,L_*,C_{\mathrm{up}},1).
\]
Their required positivity and range conditions were established above. Distinguish the restricted pointwise risk by defining
\[
R_{\mathrm{pt}}^{\mathrm{bd}}(n,\gamma,x_0)
:=
\inf_{\widehat\theta\ \mathrm{sample\text{-}measurable}}
\sup_{\substack{P\in\mathcal P_\gamma\\
P\text{ admits a causal completion with}\\
Y(0),Y(1)\in\{0,B_*\}\ \text{a.s.}}}
\mathbb E_{P^{\otimes n}}
\left|\widehat\theta(\mathcal D_n)-\mu_{1,P}(x_0)\right|.
\]
For every \(\gamma>1\) and every \(x_0\in[0,1]^d\), the cited conclusion supplies \(c_\gamma>0\) such that
\[
c_\gamma r_{\gamma,n}
\le R_{\mathrm{pt}}^{\mathrm{bd}}(n,\gamma,x_0)
\le R_{\mathrm{pt}}(n,\gamma,x_0),
\qquad n\ge1.
\]
The second inequality follows because the restricted laws form a subset of the full causal class in \cref{obj:def:risks}: for each measurable estimator, restriction decreases the worst-case risk, and taking the infimum preserves this inequality. This gives the final assertion.
% lean: CausalSmith.Stat.WeakOverlap.adaptive_minimax_rate
\end{enumerate}
\end{proof}
